\documentclass[11pt,a4paper]{article}
\usepackage[margin=2.5cm]{geometry}\usepackage{pdfpages}\usepackage{amsmath}
\usepackage{lmodern}\usepackage[T1]{fontenc}\usepackage[utf8]{inputenc}\usepackage{enumitem}\usepackage{xurl}\usepackage[colorlinks=true,urlcolor=blue!60!black,linkcolor=black,citecolor=black,bookmarks=true,bookmarksopen=true]{hyperref}\usepackage{bookmark}
\usepackage{graphicx,caption}
\begin{document}\setcounter{page}{104}
\pagestyle{plain}
\phantomsection\pdfbookmark[0]{Does everything have a centre?}{phil}
\begin{center}{\Large\bfseries Does everything have a centre? A philosophical note\par}\medskip{\large Chris Alameddin\par}\end{center}
\section*{1. The thesis}
Every physical object or system that has ever been observed has a centre. Depending on what is being described, that centre may be geometric (the centroid of a shape), dynamical (the centre of mass), or kinematic (the centre-of-momentum frame, in which the total momentum vanishes). A galaxy, a cluster, a star, a supernova remnant, a planetary system and a laboratory sample all admit one or more of these. Even an irregular or diffuse distribution of matter has a well-defined first moment,
\begin{equation}
\mathbf{x}_{\rm c}=\frac{\int\rho(\mathbf{x})\,\mathbf{x}\,d^3x}{\int\rho(\mathbf{x})\,d^3x},
\end{equation}
provided the total mass is finite. The integral fails to define a centre only when the distribution is infinite in extent or exactly uniform without bound.

The second half of the thesis is empirical. Over many discussions, the author has asked for counter-examples, including from several AI systems. Every example proposed of something \emph{without} a centre has turned out to be a mathematical abstraction, and none of them has ever been observed by any scientist as a physical object.
\section*{2. The candidate counter-examples}
The usual candidates fall into a few classes:
\begin{itemize}
\item \textbf{Infinite uniform space, an infinite plane or an infinite line.} These are idealisations. No finite measurement can establish that something extends without bound, and any observed region of them is finite and so has a centroid.
\item \textbf{The surface of a sphere, taken as an intrinsic two-dimensional manifold.} No point on the surface is special. But every physical sphere (a balloon, a planet, a soap film) is embedded in three-dimensional space and has a centre there. A centreless sphere exists only if you deliberately forget the embedding.
\item \textbf{Fractal limits.} A physical fractal-like structure (a coastline, a turbulent cascade, galaxy clustering on intermediate scales) is cut off at small and large scales. The ideal infinite-scale fractal is never observed.
\item \textbf{The homogeneous, isotropic FLRW universe.} Under the cosmological principle every point is equivalent, so there is no centre. This is the case that matters here, and it is taken up next.
\end{itemize}
In each case the absence of a centre comes from an idealising step: infinite extent, exact uniformity, or ignoring the embedding. It does not come from a measurement.
\section*{3. ``No centre'' in cosmology is an assumption}
The standard model's statement that the universe has no centre is not an observational result. It follows from the cosmological principle, which assumes spatial homogeneity and isotropy and thereby selects the FLRW metrics~\cite{rindler,ellis79}. Two statements are often merged and should be kept apart:
\begin{enumerate}
\item ``The CMB is very nearly isotropic about us.'' This is data.
\item ``Therefore there is no centre anywhere.'' This is not a deduction from (1). Isotropy about one point plus the Copernican assumption that we are not special implies homogeneity, but that assumption is a philosophical preference, not a measurement~\cite{ellis79,ellis02}.
\end{enumerate}
Spherically symmetric inhomogeneous models such as Lema\^itre--Tolman--Bondi solutions, and the earlier kinematic cosmology of Milne~\cite{milne}, show that a universe with a centre is mathematically consistent with general relativity and with many of the same observations. Ellis and collaborators have stressed that what we see on large scales depends on how local structure is averaged and fitted to a smooth model~\cite{ellis02,ellisstoeger}. Whether a centre exists is therefore an empirical question that the cosmological principle settles by assumption.
\section*{4. Connection to the CoE model}
The papers in this collection take the opposite starting point. If the universe began as an expansion from a point, a finite expansion necessarily has a centre, the centre of expansion (CoE). We sit a distance $D$ from it. Nothing in this requires local speeds above $c$: bulk recession of a distant shell faster than $c$ is expansion, not local motion.

On this view the thesis of Section~1 applies to the universe itself. A universe without a centre would be the one physical system that behaves like the mathematical abstractions of Section~2, and would be the only such system ever claimed.
\section*{5. What could test it}
A centre is not just a philosophical position. An off-centre observer should see direction-dependent signatures:
\begin{itemize}
\item a dependence of the local Hubble rate on direction relative to the CoE, measured in the papers as $A/A_{\rm all}(\theta)$, the per-angle expansion rate divided by the full-sky mean, compared with cluster-scaling anisotropy results such as those of Migkas et al.~\cite{migkas};
\item dipoles that do not line up with the kinematic CMB dipole, such as the quasar number-count dipole reported by Secrest et al.~\cite{secrest};
\item a spread in the inferred $H_0$ with direction and distance, which bears on the Hubble tension.
\end{itemize}
These signals may also have conventional explanations: survey systematics, local bulk flows, or ordinary structure. A null result at sufficient precision would push the CoE far away or remove it. A consistent, aligned set of anisotropies would be evidence that the universe, like everything else observed, has a centre.
\section*{6. Balance}
None of this proves that the universe has a centre. The argument is an inductive one: every observed system has a centre, and every centreless example is an abstraction. Induction can fail, and the universe may be unique. But the burden is symmetric. ``No centre'' should be treated as a hypothesis to be tested, on the same footing as the CoE, and not as a settled result.
\clearpage
\phantomsection\pdfbookmark[0]{References}{allrefs}
\begin{center}{\Large\bfseries References\par}\end{center}
\noindent\small References are grouped by paper. Each group keeps that paper's own numbering, so the bracketed citations in the original pages refer to the matching group below.
\subsection*{1. CoE cosmology: Apparent acceleration from an off-centre finite geometry}
\begin{description}[leftmargin=2.6em,labelwidth=2.2em,labelsep=0.4em,itemsep=1pt,parsep=0pt,font=\normalfont]
\item[{[1]}] A. G. Riess et al., Astron. J. 116, 1009 (1998). \href{https://doi.org/10.1086/300499}{\nolinkurl{doi:10.1086/300499}}
\item[{[2]}] S. Perlmutter et al., Astrophys. J. 517, 565 (1999). \href{https://doi.org/10.1086/307221}{\nolinkurl{doi:10.1086/307221}}
\item[{[3]}] G. F. R. Ellis and W. R. Stoeger, Class. Quantum Grav. 4, 1697 (1987). \href{https://doi.org/10.1088/0264-9381/4/6/025}{\nolinkurl{doi:10.1088/0264-9381/4/6/025}}
\item[{[4]}] E. Hubble, Proc. Natl. Acad. Sci. USA 15, 168 (1929). \href{https://doi.org/10.1073/pnas.15.3.168}{\nolinkurl{doi:10.1073/pnas.15.3.168}}
\item[{[5]}] K. Migkas et al., Astron. Astrophys. 636, A15 (2020). \href{https://doi.org/10.1051/0004-6361/201936602}{\nolinkurl{doi:10.1051/0004-6361/201936602}}, \href{https://arxiv.org/abs/2004.03305}{\nolinkurl{arXiv:2004.03305}}
\item[{[6]}] L. I. Sedov, Prikl. Mat. Mekh. [J. Appl. Math. Mech.] 10, 241 (1946). \textit{[no online link found]}
\item[{[7]}] G. I. Taylor, Proc. R. Soc. Lond. A 201, 159 (1950). \href{https://doi.org/10.1098/rspa.1950.0049}{\nolinkurl{doi:10.1098/rspa.1950.0049}}
\item[{[8]}] G. Lemaître, Ann. Soc. Sci. Bruxelles A53, 51 (1933); English translation: Gen. Rel. Grav. 29, 641 (1997). \href{https://doi.org/10.1023/A:1018855621348}{\nolinkurl{doi:10.1023/A:1018855621348}} (translation)
\item[{[9]}] R. C. Tolman, Proc. Natl. Acad. Sci. USA 20, 169 (1934). \href{https://doi.org/10.1073/pnas.20.3.169}{\nolinkurl{doi:10.1073/pnas.20.3.169}}
\item[{[10]}] H. Bondi, Mon. Not. R. Astron. Soc. 107, 410 (1947). \href{https://doi.org/10.1093/mnras/107.5-6.410}{\nolinkurl{doi:10.1093/mnras/107.5-6.410}}
\item[{[11]}] Planck Collaboration, Astron. Astrophys. 641, A7 (2020). \href{https://doi.org/10.1051/0004-6361/201935201}{\nolinkurl{doi:10.1051/0004-6361/201935201}}
\item[{[12]}] N. J. Secrest, S. von Hausegger, M. Rameez, R. Mohayaee, S. Sarkar, and J. Colin, Astrophys. J. Lett. 908, L51 (2021). \href{https://doi.org/10.3847/2041-8213/abdd40}{\nolinkurl{doi:10.3847/2041-8213/abdd40}}, \href{https://arxiv.org/abs/2009.14826}{\nolinkurl{arXiv:2009.14826}}
\end{description}
\noindent\textit{The supplement on the exact relativistic (LTB) calculation cites no additional references.}
\subsection*{2. Reading the CMB fingerprint of the CoE shell model}
\begin{description}[leftmargin=2.6em,labelwidth=2.2em,labelsep=0.4em,itemsep=1pt,parsep=0pt,font=\normalfont]
\item[{[1]}] J. Clampitt and B. Jain, Lensing measurements of the mass distribution in SDSS voids, MNRAS 454, 3357 (2015). \href{https://doi.org/10.1093/mnras/stv2215}{\nolinkurl{doi:10.1093/mnras/stv2215}}, \href{https://arxiv.org/abs/1404.1834}{\nolinkurl{arXiv:1404.1834}}
\item[{[2]}] J. Tinker et al., Toward a halo mass function for precision cosmology: The limits of universality, ApJ 688, 709 (2008). \href{https://doi.org/10.1086/591439}{\nolinkurl{doi:10.1086/591439}}, \href{https://arxiv.org/abs/0803.2706}{\nolinkurl{arXiv:0803.2706}}
\item[{[3]}] E. S. Rykoff et al., redMaPPer. I. algorithm and SDSS DR8 catalog, ApJ 785, 104 (2014). \href{https://doi.org/10.1088/0004-637X/785/2/104}{\nolinkurl{doi:10.1088/0004-637X/785/2/104}}, \href{https://arxiv.org/abs/1303.3562}{\nolinkurl{arXiv:1303.3562}}
\item[{[4]}] M. Simet et al., Weak lensing measurement of the mass–richness relation of SDSS redMaPPer clusters, MNRAS 466, 3103 (2017). \href{https://doi.org/10.1093/mnras/stw3250}{\nolinkurl{doi:10.1093/mnras/stw3250}}, \href{https://arxiv.org/abs/1603.06953}{\nolinkurl{arXiv:1603.06953}}
\item[{[5]}] L. J. Liivamägi, E. Tempel, and E. Saar, SDSS DR7 superclusters. The catalogues, A\&A 539, A80 (2012). \href{https://doi.org/10.1051/0004-6361/201016288}{\nolinkurl{doi:10.1051/0004-6361/201016288}}, \href{https://arxiv.org/abs/1012.1989}{\nolinkurl{arXiv:1012.1989}}
\item[{[6]}] Q. Mao et al., A cosmic void catalog of SDSS DR12 BOSS galaxies, ApJ 835, 161 (2017). \href{https://doi.org/10.3847/1538-4357/835/2/161}{\nolinkurl{doi:10.3847/1538-4357/835/2/161}}, \href{https://arxiv.org/abs/1602.02771}{\nolinkurl{arXiv:1602.02771}}
\item[{[7]}] N. Hamaus, P. M. Sutter, and B. D. Wandelt, Universal density profile for cosmic voids, PRL 112, 251302 (2014). \href{https://doi.org/10.1103/PhysRevLett.112.251302}{\nolinkurl{doi:10.1103/PhysRevLett.112.251302}}, \href{https://arxiv.org/abs/1403.5499}{\nolinkurl{arXiv:1403.5499}}
\item[{[8]}] I. Szapudi et al., Detection of a supervoid aligned with the cold spot of the cosmic microwave background, MNRAS 450, 288 (2015). \href{https://doi.org/10.1093/mnras/stv488}{\nolinkurl{doi:10.1093/mnras/stv488}}, \href{https://arxiv.org/abs/1405.1566}{\nolinkurl{arXiv:1405.1566}}
\item[{[9]}] A. Kovács and J. García-Bellido, Cosmic troublemakers: the cold spot, the eridanus supervoid, and the great walls, MNRAS 462, 1882 (2016). \href{https://doi.org/10.1093/mnras/stw1752}{\nolinkurl{doi:10.1093/mnras/stw1752}}, \href{https://arxiv.org/abs/1511.09008}{\nolinkurl{arXiv:1511.09008}}
\item[{[10]}] R. C. Keenan, A. J. Barger, and L. L. Cowie, Evidence for a $\sim$300 megaparsec scale under-density in the local galaxy distribution, ApJ 775, 62 (2013). \href{https://doi.org/10.1088/0004-637X/775/1/62}{\nolinkurl{doi:10.1088/0004-637X/775/1/62}}, \href{https://arxiv.org/abs/1304.2884}{\nolinkurl{arXiv:1304.2884}}
\item[{[11]}] J. R. Whitbourn and T. Shanks, The local hole revealed by galaxy counts and redshifts, MNRAS 437, 2146 (2014). \href{https://doi.org/10.1093/mnras/stt2024}{\nolinkurl{doi:10.1093/mnras/stt2024}}, \href{https://arxiv.org/abs/1307.4405}{\nolinkurl{arXiv:1307.4405}}
\item[{[12]}] W. D. Kenworthy, D. Scolnic, and A. Riess, The local perspective on the Hubble tension: Local structure does not impact measurement of the Hubble constant, ApJ 875, 145 (2019). \href{https://doi.org/10.3847/1538-4357/ab0ebf}{\nolinkurl{doi:10.3847/1538-4357/ab0ebf}}, \href{https://arxiv.org/abs/1901.08681}{\nolinkurl{arXiv:1901.08681}}
\item[{[13]}] J. Carrick, S. J. Turnbull, G. Lavaux, and M. J. Hudson, Cosmological parameters from the comparison of peculiar velocities with predictions from the 2M++ density field, MNRAS 450, 317 (2015). \href{https://doi.org/10.1093/mnras/stv547}{\nolinkurl{doi:10.1093/mnras/stv547}}, \href{https://arxiv.org/abs/1504.04627}{\nolinkurl{arXiv:1504.04627}}
\item[{[14]}] M. I. Scrimgeour et al., The WiggleZ dark energy survey: the transition to large-scale cosmic homogeneity, MNRAS 425, 116 (2012). \href{https://doi.org/10.1111/j.1365-2966.2012.21402.x}{\nolinkurl{doi:10.1111/j.1365-2966.2012.21402.x}}, \href{https://arxiv.org/abs/1205.6812}{\nolinkurl{arXiv:1205.6812}}
\item[{[15]}] N. J. Secrest, S. von Hausegger, M. Rameez, R. Mohayaee, S. Sarkar, and J. Colin, A test of the cosmological principle with quasars, ApJL 908, L51 (2021). \href{https://doi.org/10.3847/2041-8213/abdd40}{\nolinkurl{doi:10.3847/2041-8213/abdd40}}, \href{https://arxiv.org/abs/2009.14826}{\nolinkurl{arXiv:2009.14826}}
\item[{[16]}] K. Migkas, G. Schellenberger, T. H. Reiprich, F. Pacaud, M. E. Ramos-Ceja, and L. Lovisari, Probing cosmic isotropy with a new X-ray galaxy cluster sample through the $L_X$–$T$ scaling relation, A\&A 636, A15 (2020). \href{https://doi.org/10.1051/0004-6361/201936602}{\nolinkurl{doi:10.1051/0004-6361/201936602}}, \href{https://arxiv.org/abs/2004.03305}{\nolinkurl{arXiv:2004.03305}}
\item[{[17]}] D. J. Eisenstein and W. Hu, Baryonic features in the matter transfer function, ApJ 496, 605 (1998). \href{https://doi.org/10.1086/305424}{\nolinkurl{doi:10.1086/305424}}, \href{https://arxiv.org/abs/astro-ph/9709112}{\nolinkurl{arXiv:astro-ph/9709112}}
\item[{[18]}] J. Einasto, M. Einasto, S. Gottlöber, V. Müller, V. Saar, A. A. Starobinsky, E. Tago, D. Tucker, H. Andernach, and P. Frisch, A 120-Mpc periodicity in the three-dimensional distribution of galaxy superclusters, Nature 385, 139 (1997). \href{https://doi.org/10.1038/385139a0}{\nolinkurl{doi:10.1038/385139a0}}, \href{https://arxiv.org/abs/astro-ph/9701018}{\nolinkurl{arXiv:astro-ph/9701018}}
\item[{[19]}] J. Einasto, M. Einasto, P. Frisch, S. Gottlöber, V. Müller, V. Saar, A. A. Starobinsky, E. Tago, D. Tucker, and H. Andernach, The supercluster–void network II. An oscillating cluster correlation function, MNRAS 289, 801 (1997). \href{https://doi.org/10.1093/mnras/289.4.801}{\nolinkurl{doi:10.1093/mnras/289.4.801}}, \href{https://arxiv.org/abs/astro-ph/9704127}{\nolinkurl{arXiv:astro-ph/9704127}}
\item[{[20]}] E. Jennings, Y. Li, and W. Hu, The abundance of voids and the excursion set formalism, MNRAS 434, 2167 (2013). \href{https://doi.org/10.1093/mnras/stt1169}{\nolinkurl{doi:10.1093/mnras/stt1169}}, \href{https://arxiv.org/abs/1304.6087}{\nolinkurl{arXiv:1304.6087}}
\item[{[21]}] R. K. Sheth and R. van de Weygaert, A hierarchy of voids: much ado about nothing, MNRAS 350, 517 (2004). \href{https://doi.org/10.1111/j.1365-2966.2004.07661.x}{\nolinkurl{doi:10.1111/j.1365-2966.2004.07661.x}}, \href{https://arxiv.org/abs/astro-ph/0311260}{\nolinkurl{arXiv:astro-ph/0311260}}
\item[{[22]}] R. K. Sachs and A. M. Wolfe, Perturbations of a cosmological model and angular variations of the microwave background, ApJ 147, 73 (1967). \href{https://doi.org/10.1086/148982}{\nolinkurl{doi:10.1086/148982}}
\item[{[23]}] M. J. Rees and D. W. Sciama, Large-scale density inhomogeneities in the universe, Nature 217, 511 (1968). \href{https://doi.org/10.1038/217511a0}{\nolinkurl{doi:10.1038/217511a0}}
\end{description}
\subsection*{3. Pattern formation (v5)}
\begin{description}[leftmargin=2.6em,labelwidth=2.2em,labelsep=0.4em,itemsep=1pt,parsep=0pt,font=\normalfont]
\item[{[1]}] M. C. Cross, P. C. Hohenberg, Pattern formation outside of equilibrium, Rev. Mod. Phys. 65, 851 (1993). \href{https://doi.org/10.1103/RevModPhys.65.851}{\nolinkurl{doi:10.1103/RevModPhys.65.851}}
\item[{[2]}] D. Andelman, R. E. Rosensweig, Modulated phases: review and recent results, J. Phys. Chem. B 113, 3785 (2009). \href{https://doi.org/10.1021/jp807770n}{\nolinkurl{doi:10.1021/jp807770n}}
\item[{[3]}] J. S. Langer, Instabilities and pattern formation in crystal growth, Rev. Mod. Phys. 52, 1 (1980). \href{https://doi.org/10.1103/RevModPhys.52.1}{\nolinkurl{doi:10.1103/RevModPhys.52.1}}
\item[{[4]}] M. Seul, D. Andelman, Domain shapes and patterns: the phenomenology of modulated phases, Science 267, 476 (1995). \href{https://doi.org/10.1126/science.267.5197.476}{\nolinkurl{doi:10.1126/science.267.5197.476}}
\item[{[5]}] P. C. Hohenberg, A. P. Krekhov, An introduction to the Ginzburg–Landau theory of phase transitions and nonequilibrium patterns, Phys. Rep. 572, 1 (2015). \href{https://doi.org/10.1016/j.physrep.2015.01.001}{\nolinkurl{doi:10.1016/j.physrep.2015.01.001}}, \href{https://arxiv.org/abs/1410.7285}{\nolinkurl{arXiv:1410.7285}}
\item[{[6]}] J. Swift, P. C. Hohenberg, Hydrodynamic fluctuations at the convective instability, Phys. Rev. A 15, 319 (1977). \href{https://doi.org/10.1103/PhysRevA.15.319}{\nolinkurl{doi:10.1103/PhysRevA.15.319}}
\item[{[7]}] W. van Saarloos, Front propagation into unstable states, Phys. Rep. 386, 29 (2003). \href{https://doi.org/10.1016/j.physrep.2003.08.001}{\nolinkurl{doi:10.1016/j.physrep.2003.08.001}}, \href{https://arxiv.org/abs/cond-mat/0308540}{\nolinkurl{arXiv:cond-mat/0308540}}
\item[{[8]}] I. S. Aranson, L. S. Tsimring, Patterns and collective behavior in granular media: theoretical concepts, Rev. Mod. Phys. 78, 641 (2006). \href{https://doi.org/10.1103/RevModPhys.78.641}{\nolinkurl{doi:10.1103/RevModPhys.78.641}}, \href{https://arxiv.org/abs/cond-mat/0507419}{\nolinkurl{arXiv:cond-mat/0507419}}
\item[{[9]}] A. J. Bray, Theory of phase-ordering kinetics, Adv. Phys. 43, 357 (1994). \href{https://doi.org/10.1080/00018739400101505}{\nolinkurl{doi:10.1080/00018739400101505}}, \href{https://arxiv.org/abs/cond-mat/9501089}{\nolinkurl{arXiv:cond-mat/9501089}}
\item[{[10]}] T. D. Rossing, Chladni’s law for vibrating plates, Am. J. Phys. 50, 271 (1982). \href{https://doi.org/10.1119/1.12866}{\nolinkurl{doi:10.1119/1.12866}}
\item[{[11]}] W. S. Edwards, S. Fauve, Patterns and quasi-patterns in the Faraday experiment, J. Fluid Mech. 278, 123 (1994). \href{https://doi.org/10.1017/S0022112094003642}{\nolinkurl{doi:10.1017/S0022112094003642}}
\item[{[12]}] X. Liu, X. Wang, Polygonal patterns of Faraday water waves analogous to collective excitations in Bose–Einstein condensates, Nat. Phys. 20, 287 (2024). \href{https://doi.org/10.1038/s41567-023-02294-y}{\nolinkurl{doi:10.1038/s41567-023-02294-y}}
\item[{[13]}] S. W. Morris, E. Bodenschatz, D. S. Cannell, G. Ahlers, Spiral defect chaos in large aspect ratio Rayleigh–Bénard convection, Phys. Rev. Lett. 71, 2026 (1993). \href{https://doi.org/10.1103/PhysRevLett.71.2026}{\nolinkurl{doi:10.1103/PhysRevLett.71.2026}}
\item[{[14]}] J. R. A. Pearson, On convection cells induced by surface tension, J. Fluid Mech. 4, 489 (1958). \href{https://doi.org/10.1017/S0022112058000616}{\nolinkurl{doi:10.1017/S0022112058000616}}
\item[{[15]}] K. Eckert, M. Bestehorn, A. Thess, Square cells in surface-tension-driven Bénard convection: experiment and theory, J. Fluid Mech. 356, 155 (1998). \href{https://doi.org/10.1017/S0022112097007842}{\nolinkurl{doi:10.1017/S0022112097007842}}
\item[{[16]}] C. D. Andereck, S. S. Liu, H. L. Swinney, Flow regimes in a circular Couette system with independently rotating cylinders, J. Fluid Mech. 164, 155 (1986). \href{https://doi.org/10.1017/S0022112086002513}{\nolinkurl{doi:10.1017/S0022112086002513}}
\item[{[17]}] L. Kramer, W. Pesch, Convection instabilities in nematic liquid crystals, Ann. Rev. Fluid Mech. 27, 515 (1995). \href{https://doi.org/10.1146/annurev.fl.27.010195.002503}{\nolinkurl{doi:10.1146/annurev.fl.27.010195.002503}}
\item[{[18]}] Y. Sasaki, V. S. R. Jampani, C. Tanaka, N. Sakurai, S. Sakane, K. V. Le, F. Araoka, H. Orihara, Large-scale self-organization of reconfigurable topological defect networks in nematic liquid crystals, Nat. Commun. 7, 13238 (2016). \href{https://doi.org/10.1038/ncomms13238}{\nolinkurl{doi:10.1038/ncomms13238}}
\item[{[19]}] A. M. Turing, The chemical basis of morphogenesis, Phil. Trans. R. Soc. B 237, 37 (1952). \href{https://doi.org/10.1098/rstb.1952.0012}{\nolinkurl{doi:10.1098/rstb.1952.0012}}
\item[{[20]}] V. Castets, E. Dulos, J. Boissonade, P. De Kepper, Experimental evidence of a sustained standing Turing-type nonequilibrium chemical pattern, Phys. Rev. Lett. 64, 2953 (1990). \href{https://doi.org/10.1103/PhysRevLett.64.2953}{\nolinkurl{doi:10.1103/PhysRevLett.64.2953}}
\item[{[21]}] Q. Ouyang, H. L. Swinney, Transition from a uniform state to hexagonal and striped Turing patterns, Nature 352, 610 (1991). \href{https://doi.org/10.1038/352610a0}{\nolinkurl{doi:10.1038/352610a0}}
\item[{[22]}] I. Langmuir, Surface motion of water induced by wind, Science 87, 119 (1938). \href{https://doi.org/10.1126/science.87.2250.119}{\nolinkurl{doi:10.1126/science.87.2250.119}}
\item[{[23]}] A. D. D. Craik, S. Leibovich, A rational model for Langmuir circulations, J. Fluid Mech. 73, 401 (1976). \href{https://doi.org/10.1017/S0022112076001420}{\nolinkurl{doi:10.1017/S0022112076001420}}
\item[{[24]}] S. A. Thorpe, Langmuir circulation, Ann. Rev. Fluid Mech. 36, 55 (2004). \href{https://doi.org/10.1146/annurev.fluid.36.052203.071431}{\nolinkurl{doi:10.1146/annurev.fluid.36.052203.071431}}
\item[{[25]}] D. Etling, R. A. Brown, Roll vortices in the planetary boundary layer: a review, Boundary-Layer Meteorol. 65, 215 (1993). \href{https://doi.org/10.1007/BF00705527}{\nolinkurl{doi:10.1007/BF00705527}}
\item[{[26]}] F. Melo, P. B. Umbanhowar, H. L. Swinney, Hexagons, kinks, and disorder in oscillated granular layers, Phys. Rev. Lett. 75, 3838 (1995). \href{https://doi.org/10.1103/PhysRevLett.75.3838}{\nolinkurl{doi:10.1103/PhysRevLett.75.3838}}
\item[{[27]}] P. B. Umbanhowar, F. Melo, H. L. Swinney, Localized excitations in a vertically vibrated granular layer, Nature 382, 793 (1996). \href{https://doi.org/10.1038/382793a0}{\nolinkurl{doi:10.1038/382793a0}}
\item[{[28]}] M. D. Cowley, R. E. Rosensweig, The interfacial stability of a ferromagnetic fluid, J. Fluid Mech. 30, 671 (1967). \href{https://doi.org/10.1017/S0022112067001697}{\nolinkurl{doi:10.1017/S0022112067001697}}
\item[{[29]}] H. Knieling, R. Richter, I. Rehberg, G. Matthies, A. Lange, Growth of surface undulations at the Rosensweig instability, Phys. Rev. E 76, 066301 (2007). \href{https://doi.org/10.1103/PhysRevE.76.066301}{\nolinkurl{doi:10.1103/PhysRevE.76.066301}}, \href{https://arxiv.org/abs/0801.0140}{\nolinkurl{arXiv:0801.0140}}
\item[{[30]}] C. Kittel, Theory of the structure of ferromagnetic domains in films and small particles, Phys. Rev. 70, 965 (1946). \href{https://doi.org/10.1103/PhysRev.70.965}{\nolinkurl{doi:10.1103/PhysRev.70.965}}
\item[{[31]}] M. Kronseder, T. N. G. Meier, M. Zimmermann, M. Buchner, M. Vogel, C. H. Back, Real-time observation of domain fluctuations in a two-dimensional magnetic model system, Nat. Commun. 6, 6832 (2015). \href{https://doi.org/10.1038/ncomms7832}{\nolinkurl{doi:10.1038/ncomms7832}}
\item[{[32]}] J. A. Glazier, S. P. Gross, J. Stavans, Dynamics of two-dimensional soap froths, Phys. Rev. A 36, 306 (1987). \href{https://doi.org/10.1103/PhysRevA.36.306}{\nolinkurl{doi:10.1103/PhysRevA.36.306}}
\item[{[33]}] J. W. Miles, On the stability of heterogeneous shear flows, J. Fluid Mech. 10, 496 (1961). \href{https://doi.org/10.1017/S0022112061000305}{\nolinkurl{doi:10.1017/S0022112061000305}}
\item[{[34]}] L. N. Howard, Note on a paper of John W. Miles, J. Fluid Mech. 10, 509 (1961). \href{https://doi.org/10.1017/S0022112061000317}{\nolinkurl{doi:10.1017/S0022112061000317}}
\item[{[35]}] A. Michalke, On the inviscid instability of the hyperbolic-tangent velocity profile, J. Fluid Mech. 19, 543 (1964). \href{https://doi.org/10.1017/S0022112064000908}{\nolinkurl{doi:10.1017/S0022112064000908}}
\item[{[36]}] P. G. Saffman, G. I. Taylor, The penetration of a fluid into a porous medium or Hele-Shaw cell containing a more viscous liquid, Proc. R. Soc. Lond. A 245, 312 (1958). \href{https://doi.org/10.1098/rspa.1958.0085}{\nolinkurl{doi:10.1098/rspa.1958.0085}}
\item[{[37]}] G. M. Homsy, Viscous fingering in porous media, Ann. Rev. Fluid Mech. 19, 271 (1987). \href{https://doi.org/10.1146/annurev.fl.19.010187.001415}{\nolinkurl{doi:10.1146/annurev.fl.19.010187.001415}}
\item[{[38]}] D. H. Sharp, An overview of Rayleigh–Taylor instability, Physica D 12, 3 (1984). \href{https://doi.org/10.1016/0167-2789(84)90510-4}{\nolinkurl{doi:10.1016/0167-2789(84)90510-4}}
\item[{[39]}] W. W. Mullins, R. F. Sekerka, Stability of a planar interface during solidification of a dilute binary alloy, J. Appl. Phys. 35, 444 (1964). \href{https://doi.org/10.1063/1.1713333}{\nolinkurl{doi:10.1063/1.1713333}}
\item[{[40]}] S.-C. Huang, M. E. Glicksman, Fundamentals of dendritic solidification I. Steady-state tip growth, Acta Metall. 29, 701 (1981). \href{https://doi.org/10.1016/0001-6160(81)90115-2}{\nolinkurl{doi:10.1016/0001-6160(81)90115-2}}
\item[{[41]}] M. Matsushita, H. Fujikawa, Diffusion-limited growth in bacterial colony formation, Physica A 168, 498 (1990). \href{https://doi.org/10.1016/0378-4371(90)90402-E}{\nolinkurl{doi:10.1016/0378-4371(90)90402-E}}
\item[{[42]}] T. A. Witten, L. M. Sander, Diffusion-limited aggregation, a kinetic critical phenomenon, Phys. Rev. Lett. 47, 1400 (1981). \href{https://doi.org/10.1103/PhysRevLett.47.1400}{\nolinkurl{doi:10.1103/PhysRevLett.47.1400}}
\item[{[43]}] E. Ben-Jacob, O. Schochet, A. Tenenbaum, I. Cohen, A. Czirók, T. Vicsek, Generic modelling of cooperative growth patterns in bacterial colonies, Nature 368, 46 (1994). \href{https://doi.org/10.1038/368046a0}{\nolinkurl{doi:10.1038/368046a0}}
\item[{[44]}] A. N. Zaikin, A. M. Zhabotinsky, Concentration wave propagation in two-dimensional liquid-phase self-oscillating system, Nature 225, 535 (1970). \href{https://doi.org/10.1038/225535b0}{\nolinkurl{doi:10.1038/225535b0}}
\item[{[45]}] J. O’Donnell, S. G. Ackleson, E. R. Levine, On the spatial scales of a river plume, J. Geophys. Res. Oceans 113, C04017 (2008). \href{https://doi.org/10.1029/2007JC004440}{\nolinkurl{doi:10.1029/2007JC004440}}
\item[{[46]}] A. R. Horner-Devine, R. D. Hetland, D. G. MacDonald, Mixing and Transport in Coastal River Plumes, Ann. Rev. Fluid Mech. 47, 569 (2015). \href{https://doi.org/10.1146/annurev-fluid-010313-141408}{\nolinkurl{doi:10.1146/annurev-fluid-010313-141408}}
\item[{[47]}] J. C. McWilliams, Oceanic Frontogenesis, Ann. Rev. Mar. Sci. 13, 227 (2021). \href{https://doi.org/10.1146/annurev-marine-032320-120725}{\nolinkurl{doi:10.1146/annurev-marine-032320-120725}}
\item[{[48]}] T. Antal, M. Droz, J. Magnin, Z. Rácz, Formation of Liesegang patterns: a spinodal decomposition scenario, Phys. Rev. Lett. 83, 2880 (1999). \href{https://doi.org/10.1103/PhysRevLett.83.2880}{\nolinkurl{doi:10.1103/PhysRevLett.83.2880}}
\item[{[49]}] K. J. Måløy, J. Feder, T. Jøssang, Viscous fingering fractals in porous media, Phys. Rev. Lett. 55, 2688 (1985). \href{https://doi.org/10.1103/PhysRevLett.55.2688}{\nolinkurl{doi:10.1103/PhysRevLett.55.2688}}
\item[{[50]}] A. Libchaber, C. Laroche, S. Fauve, Period doubling cascade in mercury, a quantitative measurement, J. Physique Lett. 43, 211 (1982). \href{https://doi.org/10.1051/jphyslet:01982004307021100}{\nolinkurl{doi:10.1051/jphyslet:01982004307021100}}
\item[{[51]}] M. J. Feigenbaum, Quantitative universality for a class of nonlinear transformations, J. Stat. Phys. 19, 25 (1978). \href{https://doi.org/10.1007/BF01020332}{\nolinkurl{doi:10.1007/BF01020332}}
\item[{[52]}] J. P. Gollub, H. L. Swinney, Onset of turbulence in a rotating fluid, Phys. Rev. Lett. 35, 927 (1975). \href{https://doi.org/10.1103/PhysRevLett.35.927}{\nolinkurl{doi:10.1103/PhysRevLett.35.927}}
\item[{[53]}] E. N. Lorenz, Deterministic nonperiodic flow, J. Atmos. Sci. 20, 130 (1963). \href{https://doi.org/10.1175/1520-0469(1963)020<0130:DNF>2.0.CO;2}{\nolinkurl{doi:10.1175/1520-0469(1963)020<0130:DNF>2.0.CO;2}}
\item[{[54]}] E. T. Vishniac, The dynamic and gravitational instabilities of spherical shocks, Astrophys. J. 274, 152 (1983). \href{https://doi.org/10.1086/161433}{\nolinkurl{doi:10.1086/161433}}
\item[{[55]}] C. Michaut, C. Cavet, S. E. Bouquet, F. Roy, H. C. Nguyen, Numerical study of the Vishniac instability in supernova remnants, Astrophys. J. 759, 78 (2012). \href{https://doi.org/10.1088/0004-637X/759/2/78}{\nolinkurl{doi:10.1088/0004-637X/759/2/78}}
\item[{[56]}] R. K. Sheth, R. van de Weygaert, A hierarchy of voids: much ado about nothing, Mon. Not. R. Astron. Soc. 350, 517 (2004). \href{https://doi.org/10.1111/j.1365-2966.2004.07661.x}{\nolinkurl{doi:10.1111/j.1365-2966.2004.07661.x}}
\end{description}
\noindent\textit{Inserted Figs.~99, 100 and 104} are reproduced from Michaut et al.\ (2012), Astrophys. J. 759, 78, \href{https://doi.org/10.1088/0004-637X/759/2/78}{\nolinkurl{doi:10.1088/0004-637X/759/2/78}} (ref.~[55] above). Fig.~105 is original to this collection.
\subsection*{4. Does everything have a centre? A philosophical note}
\renewcommand{\refname}{}\vspace{-3.5em}
\begin{thebibliography}{9}\setlength{\itemsep}{1pt}
\bibitem{ellis79} G.~F.~R. Ellis, ``The homogeneity of the universe,'' \emph{Gen.\ Rel.\ Grav.} \textbf{11}, 281 (1979). \href{https://doi.org/10.1007/BF00759562}{\nolinkurl{doi:10.1007/BF00759562}}
\bibitem{ellis02} G.~F.~R. Ellis, ``Cosmology and local physics,'' \emph{New Astron.\ Rev.} \textbf{46}, 645 (2002). \href{https://doi.org/10.1016/S1387-6473(02)00234-8}{\nolinkurl{doi:10.1016/S1387-6473(02)00234-8}}
\bibitem{ellisstoeger} G.~F.~R. Ellis and W.~R. Stoeger, ``The `fitting problem' in cosmology,'' \emph{Class.\ Quantum Grav.} \textbf{4}, 1697 (1987). \href{https://doi.org/10.1088/0264-9381/4/6/025}{\nolinkurl{doi:10.1088/0264-9381/4/6/025}}
\bibitem{milne} E.~A. Milne, \emph{Relativity, Gravitation and World-Structure} (Oxford, 1935). \href{https://archive.org/details/dli.ernet.7948}{\nolinkurl{archive.org/details/dli.ernet.7948}}
\bibitem{rindler} W.~Rindler, \emph{Relativity: Special, General, and Cosmological}, 2nd ed.\ (Oxford, 2006). \href{https://doi.org/10.1093/oso/9780198567318.001.0001}{\nolinkurl{doi:10.1093/oso/9780198567318.001.0001}}
\bibitem{migkas} K.~Migkas et al., ``Probing cosmic isotropy with a new X-ray galaxy cluster sample through the $L_X$--$T$ scaling relation,'' \emph{A\&A} \textbf{636}, A15 (2020). \href{https://doi.org/10.1051/0004-6361/201936602}{\nolinkurl{doi:10.1051/0004-6361/201936602}}, \href{https://arxiv.org/abs/2004.03305}{\nolinkurl{arXiv:2004.03305}}
\bibitem{secrest} N.~J. Secrest et al., ``A test of the cosmological principle with quasars,'' \emph{ApJL} \textbf{908}, L51 (2021). \href{https://doi.org/10.3847/2041-8213/abdd40}{\nolinkurl{doi:10.3847/2041-8213/abdd40}}, \href{https://arxiv.org/abs/2009.14826}{\nolinkurl{arXiv:2009.14826}}
\end{thebibliography}

\clearpage
\setcounter{figure}{110}\captionsetup{skip=2pt,font=footnotesize,labelfont=bf}
\vspace*{\fill}\begin{center}
\includegraphics[width=\linewidth,height=0.8\textheight,keepaspectratio]{/workspace/chain/coe/figs/coe_ncp_hemisphere_path.png}
\captionof{figure}{CoE on the northern celestial hemisphere (NCP-centered): fixed ICRS direction from Galactic $(l,b)=(180^\circ,-15^\circ)$; nearest bright star $\approx$ Aldebaran ($\sim 5.3^\circ$). Uses the earlier rounded axis.}
\end{center}\vspace*{\fill}
\end{document}
